\subsection{One geometric rule and one shape parameter}

Fix a focus $F$ and a directrix $D$. A point $P$ belongs to the conic when
\[
\boxed{d(P,F)=e\,d(P,D)},
\]
where $e>0$ is the eccentricity.

The value of $e$ compares focal distance with distance from the directrix:
\[
\begin{array}{c|c}
0<e<1 & \text{ellipse},\\
e=1 & \text{parabola},\\
e>1 & \text{hyperbola}.
\end{array}
\]

\begin{center}
\begin{tikzpicture}[scale=0.82,>=stealth]
    \draw[gray,thick] (4.2,-3.5)--(4.2,3.5) node[above] {$D$};
    \coordinate (F) at (0,0);
    \coordinate (P) at (2.5,1.8);
    \coordinate (H) at (4.2,1.8);
    \draw[axis,->] (-1.0,0)--(5.4,0) node[right] {polar axis};
    \draw[blue!70!black,line width=1.1pt,->] (F)--(P)
        node[midway,above left] {$r=d(P,F)$};
    \draw[orange!85!black,line width=1.1pt] (P)--(H)
        node[midway,above] {$d(P,D)$};
    \draw (0.9,0) arc[start angle=0,end angle=35.75,radius=0.9];
    \node at (1.15,0.45) {$\theta$};
    \fill[geometricSubsetColor] (F) circle (2.2pt) node[below left] {$F$};
    \fill[geometricSubsetColor] (P) circle (2.2pt) node[above left] {$P$};
    \fill (H) circle (1.6pt);
\end{tikzpicture}
\end{center}

\begin{interpretationbox}[Why eccentricity measures shape]
Small $e$ keeps the focal distance relatively short compared with the distance
to the directrix, producing a closed curve. At $e=1$ the two distances balance.
For $e>1$ the focal distance may dominate, producing an open hyperbolic branch.
\end{interpretationbox}
